\documentclass[11pt, a4paper]{article}

\usepackage{indentfirst}
\usepackage[a4paper, margin=2.5cm]{geometry}
\usepackage{verbatim}

\setlength{\parindent}{1.25cm}

\begin{document}

\noindent


\begin{center}
\textbf{ \large {Transdisciplinaridade, Observação da Natureza e
Expressão Artística na Sociedade Maia}}\\
Carlos E. L. Pereira\textsuperscript{1} e David F. Oliveira\textsuperscript{2}

IMPA Tech - Rio de Janeiro
%Nome e endereço da Instituição

{\small N.os de Matrícula:  3262\textsuperscript{1} e 3277\textsuperscript{2}}
\end{center} 

\section{Introdução} 


A maneira como compreendemos e organizamos o conhecimento está profundamente relacionada à sociedade na qual estamos inseridos. Usualmente nos habituamos a reconhecer a ciência, a arte e a religião como campos distintos da atividade humana, cada uma com suas próprias linguagens, métodos e formas de expressão. Embora suas relações e interseções sejam frequentemente reconhecidas, essa divisão ainda molda nossa concepção de mundo. Entretanto, ao observarmos outras civilizações e contextos históricos, torna-se necessário questionar até que ponto essas classificações correspondem a formas universais de organização da experiência humana ou se são, elas próprias, construções históricas e culturais.

Essa questão torna-se ainda mais relevante quando nos voltamos às sociedades pré-colombianas das Américas. Frequentemente, os conhecimentos desses povos chegam até nós sob perspectivas problemáticas, de um lado, sua desvalorização diante daquilo que chamamos tradição científica ocidental, de outro, sua idealização como realizações extraordinárias, misteriosas, com uma suposta precisão quase inexplicável. Ambas as perspectivas podem dificultar a compreensão de como esses conhecimentos foram de fato produzidos, transmitidos e incorporados à organização da vida social. Como argumenta Johanna Broda (2012), compreender a produção de conhecimento nas sociedades pré-hispânicas exige situá-la em seu contexto histórico e cultural, considerando as relações entre observação da natureza, cosmologia, religião e organização social, sem reduzi-la aos padrões institucionais da ciência moderna ou recorrer a uma idealização acrítica dessas civilizações.

Nesse sentido, estudar as sociedades antigas do continente americano representa tanto uma oportunidade de compreender tradições intelectuais que são frequentemente deixadas de lado nas narrativas da história do conhecimento, quanto de refletir criticamente sobre nossas próprias categorias de interpretação. Entre essas sociedades, a civilização maia oferece um amplo panorama para essa investigação. Seus conhecimentos matemáticos e calendáricos, a observação sistemática dos fenômenos celestes, os registros escritos, a arquitetura monumental e as manifestações artísticas constituem diferentes evidências de uma produção intelectual complexa, profundamente relacionada à cosmologia, às práticas religiosas e à organização política. Essa relação pode ser observada, por exemplo, no Castillo de Kukulcán, em Chichén Itzá, onde a incidência da luz solar sobre a arquitetura produz, em determinadas épocas do ano, uma sequência de formas triangulares de sombra sobre a balaustrada, visualmente associadas ao corpo de uma serpente. Nesse caso, a observação dos fenômenos celestes, a organização arquitetônica e a expressão simbólica podem ser investigadas a partir de uma mesma manifestação material, cuja compreensão exige considerar também seu contexto cultural e religioso. Mais do que identificar realizações que possam ser classificadas separadamente como científicas ou artísticas segundo nossos critérios atuais, interessa compreender como essas diferentes formas de produção e expressão participavam da construção e transmissão de conhecimentos em seus contextos originais. 

A partir dessa perspectiva, este ensaio busca investigar como arte e ciência, entendidas aqui como conceitos distintos, se relacionavam na sociedade maia, considerando os fatores culturais, sociais, políticos e cosmológicos que permeavam suas formas de produção de conhecimento. Parte-se da hipótese de que determinadas manifestações materiais maias não podem ser compreendidas adequadamente quando separadas em categorias rígidas, uma vez que podiam articular entre diversas áreas do conhecimento. Isso não significa, contudo, atribuir aos maias uma compreensão da integração entre diferentes formas de conhecimento equivalente àquela que construímos atualmente.

Para desenvolver essa investigação, são analisados três casos: o Castillo de Kukulcán, em Chichén Itzá, pelas relações entre arquitetura, fenômenos solares e práticas rituais; as tabelas astronômicas do Códice de Dresden, pela sistematização de conhecimentos calendáricos e pela organização visual dos registros numéricos; e os murais de Bonampak, por suas dimensões narrativa, simbólica e política. A análise distingue evidências documentadas, interpretações dos pesquisadores e hipóteses ainda debatidas, buscando compreender tanto a produção de conhecimento na sociedade maia quanto os limites das classificações modernas que utilizamos para interpretar as diferentes manifestações da experiência humana.

\end{document}
